\paragraph{Validation scope.}
The 19 focused tests pass in 20.486 seconds. New tests compare 400 small
multigraphs with all their spanning trees, including loops, parallel edges,
unreachable vertices, zero costs and 512-bit costs. Genuine source tests
replay both selectors, preserve the lower charts, compare retained costs,
exercise twelve relabellings, and check invariance of the top trace and raw
costs under scaling by $2^{16384}$. The two new positive cuts replay with
both boundary backends and with the producer helpers disabled. The
abstract counterexample above is also a regression test.

\paragraph{Paired full-query measurements.}
The retained queue reducer and whole-cut wrapper are unmodified modules
from commit \texttt{97e16378b19f}. The baseline wrapper calls that retained
reducer; all other dependencies and every verifier are shared and unchanged.
Each successful measurement includes full certificate production and
independent replay. Five randomly ordered queue/weighted pairs and three
queue/queue controls were run per source, before the full test suite or
source audit began. The table reports marginal median times and the median
of paired queue/weighted ratios, which are different statistics.

\begin{center}
\small
\begin{tabular}{lrrrr}
\toprule
Source & Queue (ms) & Weighted (ms) & Paired gain & A/A \\
\midrule
No-top pilot & 40.88 & 39.98 & 1.026 & 0.996 \\
Interior, mixed & 179.43 & 181.74 & 0.981 & 0.956 \\
Figure-eight, mixed & 6239.78 & 6245.84 & 1.000 & 1.012 \\
Trefoil pilot & 12.57 & 12.03 & 1.010 & 0.995 \\
Nested 12, empty & 41.62 & 41.85 & 0.999 & 1.003 \\
Nested 12, boundary & 60.25 & 60.19 & 1.000 & 0.995 \\
Loop, empty & 1.09 & 5.28 & --- & 1.005 \\
Loop, boundary & 1.22 & 11.48 & --- & 0.976 \\
\bottomrule
\end{tabular}
\end{center}


The first six cases complete in both modes. Their paired ratios span
0.981--1.026, while their A/A medians span 0.956--1.012; these measurements
show no broad elapsed-time advantage. The no-top pilot does not execute
forest selection and is a control for this variation. The last two queue
calls return inconclusive at the core precheck, while weighted calls
produce complete independently verified censuses. Their times therefore
measure different outcomes; no speedup ratio is reported. The weighted
positive certificates have 8,963 and 15,371 serialized bytes, respectively.
The default change is justified by the exact retained-incidence objective
and the measured additional coverage, not a claimed uniform speedup.

\paragraph{Testing the next restricted forest theorem.}
A separate probe runs report 99's polynomial solver on the 120 original
core matrices, before any additional bundle-interface rows. Its eighteen
standalone tests pass from an isolated copy in 0.004 seconds. All completed
matrix answers independently replay: 83 have a forest and complete raw
peeling, two have negative forest certificates, and 35 lie outside its
at-most-two-rows-per-active-column hypothesis. Among the 24 cores with
three-handles, these counts are nine, two and thirteen. No bare core that
the new weighted selector fails to peel gains completion from this
restricted solver on this bank. This is useful scope evidence against
assuming that a general finite matroid implementation immediately adds
native coverage. The solver remains a research dependency, and its
negative answers concern the fixed matrix/forest model only.

All source pins, retained baseline modules, raw paired samples, audit
records and test logs are in \path{synthesis/data/weighted-forest-*}.
The source-to-pattern transport, complete useful-surface discovery,
hierarchy depth, restart accounting and native orbit-trace bounds needed
for general quasi-polynomial recognition remain open.
